% ============================================================
%  ACL 2023 format
%  Style files needed in the SAME folder as this .tex file:
%    acl.sty        (from acl2023 package)
%    acl_natbib.bst (from acl2023 package — already in your acl2023/ folder)
%
%  Copy acl.sty into this folder, then compile with:
%    pdflatex report.tex
%    pdflatex report.tex
%  (References are inline — no bibtex step needed)
% ============================================================

\documentclass[11pt]{article}
\usepackage[hyperref]{acl}
\setcitestyle{numbers}   % use [1][2]... style; avoids natbib author-year error
\usepackage{times}
\usepackage{latexsym}
\usepackage[T1]{fontenc}
\usepackage[utf8]{inputenc}
\usepackage{microtype}
\usepackage{inconsolata}
\usepackage{booktabs}
\usepackage{multirow}
\usepackage{amsmath}
\usepackage{xurl}

\title{Aspect-Aware Sentiment Mining and LLM-Driven Ad Generation\\
       on Amazon Musical Instrument Reviews}

\author{
  Jiayi Song \quad  Kieu Huyen Tran Tran\\[2pt]
  Bocconi University, Milan \\
}

\begin{document}
\maketitle

% ------------------------------------------------------------------
\begin{abstract}
We present a two-stage NLP pipeline applied to Amazon Musical
Instrument reviews.
\textit{Stage~1} combines LDA topic modeling with six sentiment
analysis approaches (across eight models) to extract feature-level
signals: which product attributes customers praise or criticize.
\textit{Stage~2} uses these signals to generate targeted
advertisements with a large language model, contrasting a
\textit{strength-amplifying} strategy (highlighting praised features)
against a \textit{problem-solving} strategy (addressing criticisms);
an independent LLM judge evaluates the ads via pairwise and pointwise
protocols.
On the guitar corpus (134{,}068 reviews), fine-tuned RoBERTa
\cite{liu2019roberta} achieves the highest document-level Macro~F1
(0.944), while BART zero-shot NLI \cite{lewis2020bart} leads among
training-free methods (0.911).
Topic-level analysis identifies shipping damage as the most
criticized feature (82.6\% negative) and acoustic tone the most
praised (96.2\% positive).
In ad evaluation, the problem-solving strategy wins 6/9 pairwise
matchups and leads on all four pointwise dimensions, consistent
with Prospect Theory \cite{kahneman1979prospect}.
We replicate the full pipeline on \textit{drums} and
\textit{keyboards}: Strategy~B again wins both subcategories
(5/8 pairwise wins, leading on every pointwise dimension), confirming
that the problem-solving advantage generalizes beyond guitars.
\end{abstract}

% ------------------------------------------------------------------
\section{Introduction}

Customer reviews on Amazon contain detailed, product-specific
opinions, but most marketing pipelines use them only as star-rating
averages.
We ask two linked questions: (1)~\textit{Which features of musical
instruments do customers consistently praise or criticize?} and
(2)~\textit{Can feature-level sentiment drive more persuasive
advertising copy?}

Our pipeline operates on the Amazon Reviews 2023 Musical Instruments
dataset \cite{hou2024bridging}.
We first apply Latent Dirichlet Allocation (LDA) \cite{blei2003lda}
to discover product feature topics, then benchmark eight sentiment
classifiers to validate reliable sentiment detection in this domain.
We then perform topic-level sentiment analysis to rank features by
customer satisfaction.
In Stage~2 we generate pairs of ads per product, one amplifying
praised features (\textit{Strategy~A}) and one addressing criticisms
(\textit{Strategy~B}), and evaluate them with an independent LLM
judge.
The guitar subcategory serves as our primary corpus; we then
replicate the entire pipeline on drums and keyboards
(Section~\ref{sec:rep}) to test cross-instrument generalisation.

% ------------------------------------------------------------------
\section{Dataset and EDA}
\label{sec:data}

\paragraph{Dataset.}
We use Amazon Reviews 2023 Musical Instruments
\cite{hou2024bridging}, joining reviews with product metadata on
\texttt{parent\_asin} (100\% join rate).
The full dataset contains 3{,}017{,}439 reviews across 14
subcategories with $\geq$5{,}000 reviews.
We select the \textbf{Guitars} subcategory (134{,}068 reviews)
based on a multi-criteria scorecard balancing corpus size,
review length, and sentiment balance, where guitars score second
overall (5.74/10) with a favorable size and topic diversity.

\paragraph{EDA key findings.}
The rating distribution is heavily skewed (about 6.1$\times$ more
positive than negative reviews), so we use Macro~F1 as the primary
metric since accuracy alone is misleading.
Most reviews (92.1\%) are verified, and verified vs.\ non-verified
samples show no systematic difference large enough to warrant
removal.
Reviews are short (median 23 words), and roughly a quarter contain
fewer than 10 words, motivating our preprocessing filter.

\paragraph{Binary labels and split.}
Ratings 4-5 map to \textit{positive}, ratings 1-2 to
\textit{negative}, and rating~3 is dropped as ambiguous.
We use a stratified 80/20 split (test set: 24{,}815 reviews).
Full EDA tables and the rating-3-free split footnote are in
Appendix~\ref{app:repro}.

% ------------------------------------------------------------------
\section{Topic Modeling}
\label{sec:topics}

We explored BERTopic \cite{grootendorst2022bertopic}
(all-mpnet-base-v2 + UMAP + HDBSCAN) but obtained $\approx$50\%
outlier reviews regardless of parameter settings, consistent with
known limitations of density clustering on short, mixed-register
text.
We therefore proceeded with LDA.

\paragraph{Choice of K.}
We grid-searched the number of topics over $K \in \{5, 10, \ldots, 25\}$
using C$_v$ coherence \cite{roder2015coherence}.
Coherence increases monotonically with $K$ in our range and peaks
at $K{=}25$ (coherence~=~0.541), which we adopt as our base model.
The absolute coherence ceiling near 0.54 is a known property of short
Amazon review text rather than a modeling failure.
Other LDA hyperparameters ($\alpha$, $\eta$, passes) are tuned at the
same time but reported in Appendix~\ref{app:repro}.

\paragraph{Topic labeling and assignment.}
From the 25 raw topics we manually retain the 13 interpretable ones
and merge two accessory-related topics, giving \textbf{14 final
topics} (e.g., \textit{Tuning stability}, \textit{String quality},
\textit{Shipping damage}, \textit{Beginner learning}).
Each review is then assigned to its highest-probability topic within
the 14 above a small probability threshold, which yields
\textbf{100{,}248 assigned reviews (74.8\%)}.

% ------------------------------------------------------------------
\section{Sentiment Analysis}
\label{sec:sentiment}

\subsection{Document-Level Classification}
\label{sec:doclevel}

We benchmark six approaches (eight models) on the same test set,
using \textbf{Macro~F1} as the primary metric.

\paragraph{Classical baselines.}
TF-IDF + Logistic Regression (LR) with bigrams achieves Macro~F1~=~0.893 (tuning details in Appendix~\ref{app:repro}).
Sentence Transformer embeddings
(\href{https://huggingface.co/sentence-transformers/all-mpnet-base-v2}{\texttt{all-mpnet-base-v2}})
\cite{reimers2019sbert} with LR improve marginally to 0.898.

\paragraph{Pre-trained models (zero-shot).} \href{https://huggingface.co/nlptown/bert-base-multilingual-uncased-sentiment}{\texttt{nlptown-BERT}}, binarized by summing $P(\text{4}{\star})+P(\text{5}{\star})$,
achieves 0.901.
\href{https://huggingface.co/cardiffnlp/twitter-roberta-base-sentiment-latest}{\texttt{CardiffNLP-RoBERTa}}
underperforms (0.875), likely due to Twitter domain mismatch.
\href{https://huggingface.co/facebook/bart-large-mnli}{\texttt{BART-MNLI}}
\cite{lewis2020bart} with NLI hypothesis \textit{``The sentiment of
this review is \{positive / negative\}''} achieves \textbf{0.911},
best among training-free methods.
\href{https://huggingface.co/cross-encoder/nli-deberta-v3-large}{\texttt{DeBERTa-NLI}}
reaches 0.901.

\paragraph{Fine-tuned and few-shot.}
Fine-tuned \href{https://huggingface.co/FacebookAI/roberta-base}{\texttt{roberta-base}}
\cite{liu2019roberta} achieves \textbf{Macro~F1~=~0.944}, the
highest overall.
\href{https://huggingface.co/mistralai/Mistral-7B-Instruct-v0.3}{\texttt{Mistral-7B-Instruct-v0.3}}
\cite{jiang2023mistral} with 5-shot in-context learning
\cite{brown2020fewshot} reaches 0.924, only 0.020 below fine-tuned
RoBERTa despite using no labeled training data.
Training and prompt details are in Appendix~\ref{app:repro}.

Full results are in Table~\ref{tab:sentiment}.
Methods cluster into three tiers:
(i) \textit{top} (Macro~F1 $\geq$0.92): fine-tuned RoBERTa and
Mistral-7B;
(ii) \textit{upper} (0.90-0.92): BART, DeBERTa, nlptown, ST+LR;
(iii) \textit{lower} ($<$0.90): TF-IDF+LR and CardiffNLP.
Across topics, \textit{Tuning stability} is consistently easiest
and \textit{Pickups} hardest for all methods.

\begin{table}[t]
\centering
\small
\setlength{\tabcolsep}{4pt}
\caption{Document-level sentiment on the guitar test set.
Macro~F1 is the primary metric; Neg~F1 reflects minority-class
performance.}
\label{tab:sentiment}
\begin{tabular}{lccc}
\toprule
\textbf{Method} & \textbf{Acc} & \textbf{Mac F1} & \textbf{Neg F1} \\
\midrule
TF-IDF + LR            & 0.948 & 0.893 & 0.816 \\
SentTransf.\ + LR      & 0.950 & 0.898 & 0.825 \\
nlptown BERT           & 0.944 & 0.901 & 0.835 \\
CardiffNLP RoBERTa     & 0.933 & 0.875 & 0.790 \\
BART zero-shot         & 0.952 & 0.911 & 0.851 \\
DeBERTa zero-shot      & 0.950 & 0.901 & 0.831 \\
Fine-tuned RoBERTa     & \textbf{0.972} & \textbf{0.944} & \textbf{0.905} \\
Mistral-7B 5-shot      & 0.959 & 0.924 & 0.874 \\
\bottomrule
\end{tabular}
\end{table}

\subsection{Error Analysis}
\label{sec:error}

We inspect the 701 errors (2.8\%) made by fine-tuned RoBERTa on the
guitar test set to understand systematic failure modes.
False positives (358) and false negatives (343) are nearly balanced,
indicating no directional bias.
Notably, 84\% of errors are \textit{high-confidence} mistakes
(predicted probability $>$0.90 or $<$0.10), showing that failures
are not concentrated at the decision boundary but reflect
systematic, recurring patterns in the data.

\paragraph{Per-topic error rates.}
Error rates vary substantially by topic, with shipping- and
service-related topics hardest and acoustic-tone topics easiest
(full breakdown in Appendix~\ref{app:repro}).
A useful diagnostic is \textit{Pickups}: its error rate is low
(1.5\%) but its Macro~F1 is the lowest of any topic (0.855),
because only 38 of 1{,}477 Pickups test reviews are negative, so a
handful of minority-class mistakes depresses Macro~F1 without
moving accuracy.

\paragraph{Label noise: rating-text mismatch.}
Many errors come from reviews whose text contradicts the star
rating, e.g.\ a 4-star review whose text says \textit{``strings come
out of tune so easily, also makes a weird noise on the amp?''}.
The model correctly infers negative sentiment but is penalised by
the positive label, exposing a fundamental limit of rating-derived
labels for fine-grained sentiment.

\paragraph{Mixed sentiment.}
The second pattern involves reviews containing genuine positive and
negative signals for \textit{different} aspects, e.g.\ a 5-star
review reading \textit{``Everything is amazing with this guitar!!\
But, the amp is absolute junk!''} where the negative half drives the
model's prediction.
This is the primary motivation for topic-level sentiment analysis
(Section~\ref{sec:toplevel}): document-level labels cannot capture
within-review variation across aspects.

\subsection{Topic-Level Sentiment Analysis}
\label{sec:toplevel}

Document-level labels cannot capture within-review sentiment
variation, since a reviewer may praise acoustic tone while
criticizing string quality.
We therefore employ two approaches to topic-level sentiment.

\paragraph{Targeted sentiment via LLM (primary).}
Mistral-7B with 5-shot prompting classifies sentiment
\textit{specifically toward each review's dominant LDA topic}
for LDA-assigned reviews
(excluding 414 reviews with text length $\leq$ 10
characters).
Because LDA has already established topic relevance, no confidence
threshold is needed.
Results are in Table~\ref{tab:topics}.

\paragraph{ABSA via zero-shot NLI (secondary).}
Using BART-large-MNLI we classify every review against every one of
the 14 topic hypotheses, keeping only assignments above a confidence
threshold (see Appendix~\ref{app:repro}).
All 14 topics score net positive under ABSA, compressed into the
narrow band of 65.6\%-80.6\%, which we attribute to the
dataset's strong positive skew (80\% positive reviews) pushing
absolute scores toward the positive pole.
By contrast, the LLM approach yields a far wider spread
(17.4\%-96.2\%), with five topics scoring net negative.
The two approaches agree only at the extremes (acoustic tone ranks
first in both; shipping damage ranks last in both) but diverge
substantially in the middle, for example Fret/neck setup
(76.1\% ABSA vs.\ 43.5\% LLM) and String quality
(74.6\% vs.\ 27.6\%).
We read this as a methodological artifact: ABSA scores every review
against every topic regardless of relevance, while the LLM approach
operates only on LDA-confirmed assignments.
We therefore use LLM-based targeted sentiment as the primary source
for downstream analysis, and treat the ABSA agreement at the extremes
as a sanity check only.

\begin{table}[t]
\centering
\small
\setlength{\tabcolsep}{4pt}
\caption{Topic-level sentiment distribution (guitars).
Source: Mistral-7B targeted sentiment on LDA-assigned reviews.}
\label{tab:topics}
\begin{tabular}{lrcc}
\toprule
\textbf{Topic} & \textbf{N} & \textbf{Pos\%} & \textbf{Neg\%} \\
\midrule
Acoustic tone            & 3{,}649  & 96.2 &  3.8 \\
Beginner learning        & 12{,}509 & 91.3 &  8.7 \\
Pickups                  & 5{,}162  & 88.8 & 11.2 \\
Guitar size              & 4{,}594  & 82.0 & 18.0 \\
Playability / chords     & 4{,}437  & 81.4 & 18.6 \\
Setup / action           & 3{,}844  & 79.9 & 20.1 \\
Visual appearance        & 3{,}342  & 68.9 & 31.1 \\
Accessories              & 7{,}147  & 67.7 & 32.3 \\
Electronics / controls   & 1{,}438  & 62.8 & 37.2 \\
Customer service / ret.\ & 5{,}466  & 47.3 & 52.7 \\
Fret / neck setup        & 10{,}203 & 43.5 & 56.5 \\
Tuning stability         & 5{,}053  & 33.2 & 66.8 \\
String quality           & 4{,}820  & 27.6 & 72.4 \\
Shipping damage          & 4{,}497  & 17.4 & 82.6 \\
\bottomrule
\end{tabular}
\end{table}

% ------------------------------------------------------------------
\section{Ad Generation and Evaluation}
\label{sec:ads}

\subsection{Ad Generation}
\label{sec:adgen}

\paragraph{Product and topic selection.}
Rather than fixing topics at market level, we take a product-first
approach so that both strategies receive equally rich source
material.
For each product we keep only topics with at least 5 positive and 5
negative reviews, rank products by how many such topics they have,
deduplicate by brand, and pick the top three; within each product we
then pick the three topics closest to a 50/50 positive/negative
split.
The full eligibility numbers and the scoring formula are in
Appendix~\ref{app:select}.

The three selected guitars are entry-level products: the Davison
Full Size Electric Guitar, the Directly Cheap Acoustic Pack, and
the Crescent MG38 Acoustic Starter Pack.
\textit{String quality} appears in all three, which lets us compare
strategies on the same topic across different products.

\paragraph{Two-strategy generation.}
GPT-5.4 generates all ads in a two-step pipeline inspired by
Chain-of-Thought prompting \cite{wei2022chain}: Step~1 extracts
the most compelling genuine praise (\textit{insight\_a}) and the
most common real concern (\textit{insight\_b}); Step~2 generates
both strategies in a single call with both insights as context,
preventing strategy convergence.

\textit{Strategy~A} (strength-amplifying) adopts a celebratory,
aspirational tone grounded in \textit{insight\_a}, mapping to
Rossiter-Percy transformational motivation
\cite{rossiter1991grid}.
\textit{Strategy~B} (problem-solving) uses a reassurance framing
with the signature \textit{``Not every [product] does X, but this
one does''}, mapping to informational motivation.
Total: 3 products $\times$ 3 topics $\times$ 2 strategies = 18 ads.

\subsection{LLM-as-Judge Evaluation}
\label{sec:judge}

DeepSeek-V4-Flash serves as judge.
Using a different company and architecture from the generator
strongly mitigates self-enhancement bias \cite{zheng2023judging}.
Strategy labels are withheld from the judge, and ad order is
randomized to control for position bias.

\paragraph{Pairwise evaluation (primary).}
Each of the 9 product$\times$topic pairs is evaluated twice
(A-first and B-first) and results averaged.
Strategy~B wins 6/9 matchups (66.7\%); 2 ties; 1 Strategy~A win.

\paragraph{Pointwise evaluation (supplementary).}
Each ad is scored independently on four dimensions (1-5):
Emotional Resonance (ELM peripheral route), Argument Strength
(ELM central route) \cite{petty1986elm}, Specificity, and Purchase
Motivation (AIDA).
Chain-of-thought reasoning is required before scoring.

\paragraph{Results.}
Table~\ref{tab:ads} summarises both protocols.
Both methods consistently favor Strategy~B (composite 3.17 vs.\
2.69, $\Delta{=}0.48$).

\begin{table}[t]
\centering
\small
\setlength{\tabcolsep}{4pt}
\caption{Ad evaluation results. Pairwise: win counts (9 matchups).
Pointwise: mean scores (1-5 scale).}
\label{tab:ads}
\begin{tabular}{lcc}
\toprule
& \textbf{Strat.\ A} & \textbf{Strat.\ B} \\
\midrule
\multicolumn{3}{l}{\textit{Pairwise (primary)}} \\
\quad Wins (Ties: 2)  & 1 / 9  & \textbf{6 / 9} \\[3pt]
\multicolumn{3}{l}{\textit{Pointwise (supplementary)}} \\
\quad Emot.\ Resonance   & 2.78 & \textbf{3.33} \\
\quad Argument Strength  & 2.56 & \textbf{3.00} \\
\quad Specificity        & 2.56 & \textbf{3.00} \\
\quad Purchase Motiv.\   & 2.89 & \textbf{3.33} \\
\quad \textbf{Composite} & 2.69 & \textbf{3.17} \\
\bottomrule
\end{tabular}
\end{table}

\paragraph{Discussion.}
Strategy~B winning on Emotional Resonance goes against ELM
\cite{petty1986elm}, which predicts the aspirational
(transformational) strategy should have the emotional edge.
A possible explanation is loss aversion \cite{kahneman1979prospect}:
the ``Not every X\ldots'' contrast frame triggers fear of a bad
outcome, which produces a stronger emotional response than direct
aspiration.
The single Strategy~A win (Crescent MG38 $\times$ String quality)
fits the mixed gain/loss framing evidence in the literature
\cite{rothman1997shaping}: gain framing can outperform problem
framing when positive aspiration is more credible than implied
problem acknowledgment for a specific product-topic pair.

% ------------------------------------------------------------------
\section{Replication: Drums and Keyboards}
\label{sec:rep}

\begin{table}[t]
\centering
\small
\setlength{\tabcolsep}{4pt}
\caption{Cross-subcategory ad evaluation summary.
Pairwise: A/B/Tie counts.
Pointwise: composite mean (1-5).}
\label{tab:rep}
\begin{tabular}{lccc}
\toprule
\textbf{Subcategory} & \textbf{Pairwise} & \textbf{A} & \textbf{B} \\
\midrule
Guitars     & 1\,/\,6\,/\,2 & 2.69 & \textbf{3.17} \\
Drums       & 1\,/\,5\,/\,2 & 2.81 & \textbf{3.56} \\
Keyboards   & 0\,/\,5\,/\,3 & 2.56 & \textbf{3.12} \\
\bottomrule
\end{tabular}
\end{table}

To assess generalisability we apply the identical pipeline to the
\textit{Drums \& Percussion} and \textit{Keyboards \& MIDI}
subcategories of the same dataset, whose feature landscapes
(\textit{drums}: percussion hardware and meditation chimes and \textit{keyboards}: key action and
built-in sounds) differ markedly from guitars.

\paragraph{LDA and topic-level sentiment.}
The same grid search yields $K{=}10$ for drums (coherence~=~0.656)
and $K{=}25$ for keyboards (coherence~=~0.576); the lower drums
optimum reflects a smaller, more thematically concentrated corpus.
After manual labeling we retain \textbf{8 topics} for drums (e.g.,
\textit{Drum hardware/shells}, \textit{Product defects/returns}) and
\textbf{14 topics} for keyboards (e.g., \textit{Key feel/weighted
keys}, \textit{Returns/broken keys}).
Each LDA-assigned review is scored with the fine-tuned RoBERTa
classifier from Section~\ref{sec:doclevel}.
The most criticised topic is \textit{Product defects/returns} for
drums (55.4\% negative) and \textit{Returns/broken keys} for
keyboards (68.6\% negative).
Functional-quality issues thus dominate the negative pole in every
subcategory, mirroring guitars (full distributions in
Appendix~\ref{app:rep_topics}).

\paragraph{Ad generation and evaluation.}
We re-run the product-first selection and generate ads for the top
three products per subcategory (8 product$\times$topic pairs each, as product per subcategory has only 2 topics meeting the review count threshold).
Strategy~B wins both replications and leads on all pointwise
dimensions (Table~\ref{tab:rep}), with margins larger than for
guitars ($\Delta{=}0.75$ drums, $\Delta{=}0.56$ keyboards vs.\
$\Delta{=}0.48$).
The lone drums Strategy~A win (Woodstock Chimes $\times$ Cymbals)
parallels Crescent $\times$ String quality on guitars.

% ------------------------------------------------------------------
\section{Conclusion}
\label{sec:conclusion}

Across guitars, drums, and keyboards, problem-solving ads consistently outperformed strength-amplifying ones, with functional-quality issues proving the most reliable hook for persuasion, though future work should validate these findings with larger samples and human evaluators.

% ------------------------------------------------------------------
\section*{Acknowledgements}
We thank the Bocconi University NLP course staff for guidance
and Bocconi University for HPC access.

% ------------------------------------------------------------------
% REFERENCES — inline, no bibtex needed
% ------------------------------------------------------------------
\begin{thebibliography}{10}

\bibitem{hou2024bridging}
Yupeng Hou, Jiacheng Li, Zhankui He, An Yan, Xiusi Chen, and
Julian McAuley.
\newblock Bridging Language and Items for Retrieval and Recommendation.
\newblock \textit{arXiv preprint arXiv:2403.03952}, 2024.

\bibitem{blei2003lda}
David M. Blei, Andrew Y. Ng, and Michael I. Jordan.
\newblock Latent {D}irichlet Allocation.
\newblock \textit{Journal of Machine Learning Research}, 3:993--1022, 2003.

\bibitem{jiang2023mistral}
Albert Q. Jiang, Alexandre Sablayrolles, Arthur Mensch, Chris Bamford,
Devendra Singh Chaplot, Diego de las Casas, Florian Bressand,
Gianna Lengyel, Guillaume Lample, Lucile Saulnier, et al.
\newblock Mistral 7B.
\newblock \textit{arXiv preprint arXiv:2310.06825}, 2023.

\bibitem{wei2022chain}
Jason Wei, Xuezhi Wang, Dale Schuurmans, Maarten Bosma, Fei Xia,
Ed Chi, Quoc V. Le, and Denny Zhou.
\newblock Chain-of-{T}hought Prompting Elicits Reasoning in Large Language
  Models.
\newblock In \textit{NeurIPS}, 2022.

\bibitem{zheng2023judging}
Lianmin Zheng, Wei-Lin Chiang, Ying Sheng, Siyuan Zhuang,
Zhanghao Wu, Yonghao Zhuang, Zi Lin, Zhuohan Li, Dacheng Li,
Eric P. Xing, Hao Zhang, Joseph E. Gonzalez, and Ion Stoica.
\newblock Judging {LLM}-as-a-Judge with {MT}-Bench and Chatbot Arena.
\newblock In \textit{NeurIPS}, 2023.

\bibitem{kahneman1979prospect}
Daniel Kahneman and Amos Tversky.
\newblock Prospect Theory: An Analysis of Decision under Risk.
\newblock \textit{Econometrica}, 47(2):263--291, 1979.

\bibitem{rothman1997shaping}
Alexander J. Rothman and Peter Salovey.
\newblock Shaping Perceptions to Motivate Healthy Behavior: The Role of
  Message Framing.
\newblock \textit{Psychological Bulletin}, 121(1):3--19, 1997.

\bibitem{lewis2020bart}
Mike Lewis, Yinhan Liu, Naman Goyal, Marzieh Ghazvininejad,
Abdelrahman Mohamed, Omer Levy, Veselin Stoyanov, and Luke Zettlemoyer.
\newblock {BART}: Denoising Sequence-to-Sequence Pre-training for Natural
  Language Generation, Translation, and Comprehension.
\newblock In \textit{ACL}, 2020.

\bibitem{reimers2019sbert}
Nils Reimers and Iryna Gurevych.
\newblock Sentence-{BERT}: Sentence Embeddings using {S}iamese
  {BERT}-Networks.
\newblock In \textit{EMNLP-IJCNLP}, 2019.

\bibitem{liu2019roberta}
Yinhan Liu, Myle Ott, Naman Goyal, Jingfei Du, Mandar Joshi,
Danqi Chen, Omer Levy, Mike Lewis, Luke Zettlemoyer, and Veselin Stoyanov.
\newblock {R}o{BERT}a: A Robustly Optimized {BERT} Pretraining Approach.
\newblock \textit{arXiv preprint arXiv:1907.11692}, 2019.

\bibitem{roder2015coherence}
Michael R{\"o}der, Andreas Both, and Alexander Hinneburg.
\newblock Exploring the Space of Topic Coherence Measures.
\newblock In \textit{Proceedings of the 8th ACM International Conference on
  Web Search and Data Mining (WSDM)}, pages 399--408, 2015.

\bibitem{grootendorst2022bertopic}
Maarten Grootendorst.
\newblock {BERT}opic: Neural Topic Modeling with a Class-based {TF-IDF}
  Procedure.
\newblock \textit{arXiv preprint arXiv:2203.05794}, 2022.

\bibitem{petty1986elm}
Richard E. Petty and John T. Cacioppo.
\newblock \textit{Communication and Persuasion: Central and Peripheral
  Routes to Attitude Change}.
\newblock Springer-Verlag, New York, 1986.

\bibitem{rossiter1991grid}
John R. Rossiter, Larry Percy, and Robert J. Donovan.
\newblock A Better Advertising Planning Grid.
\newblock \textit{Journal of Advertising Research}, 31(5):11--21, 1991.

\bibitem{brown2020fewshot}
Tom B. Brown, Benjamin Mann, Nick Ryder, Melanie Subbiah,
Jared Kaplan, Prafulla Dhariwal, Arvind Neelakantan, Pranav Shyam,
Girish Sastry, Amanda Askell, et al.
\newblock Language Models are Few-Shot Learners.
\newblock In \textit{NeurIPS}, 2020.

\end{thebibliography}

% ------------------------------------------------------------------
% APPENDIX (does not count toward 4-page limit)
% ------------------------------------------------------------------
\appendix

\section{Per-Topic Breakdown (All Methods)}
\label{app:pertopic}

Table~\ref{tab:pertopic} reports per-topic Macro~F1 for all
document-level methods. \textit{Pickups} is consistently hardest;
\textit{Tuning stability} consistently easiest.

\begin{table*}[t]
\centering
\small
\setlength{\tabcolsep}{3.5pt}
\caption{Per-topic Macro~F1, guitar test set.}
\label{tab:pertopic}
\begin{tabular}{lcccccccc}
\toprule
\textbf{Topic}
  & \textbf{TF-IDF}
  & \textbf{ST+LR}
  & \textbf{nlptown}
  & \textbf{Cardiff}
  & \textbf{BART}
  & \textbf{DeBERTa}
  & \textbf{FT-Rob.}
  & \textbf{Mistral} \\
\midrule
Accessories              & 0.803 & 0.838 & 0.862 & 0.838 & 0.848 & 0.853 & 0.914 & 0.885 \\
Acoustic tone            & 0.821 & 0.872 & 0.808 & 0.785 & 0.864 & 0.810 & 0.876 & 0.927 \\
Beginner learning        & 0.861 & 0.887 & 0.902 & 0.859 & 0.887 & 0.897 & 0.937 & 0.908 \\
Customer service / ret.\ & 0.868 & 0.873 & 0.876 & 0.848 & 0.913 & 0.894 & 0.930 & 0.917 \\
Electronics / controls   & 0.816 & 0.837 & 0.879 & 0.794 & 0.820 & 0.821 & 0.908 & 0.868 \\
Fret / neck setup        & 0.884 & 0.880 & 0.887 & 0.844 & 0.908 & 0.892 & 0.939 & 0.906 \\
Guitar size              & 0.760 & 0.825 & 0.886 & 0.790 & 0.836 & 0.814 & 0.886 & 0.896 \\
Pickups                  & 0.824 & 0.763 & 0.711 & 0.692 & 0.776 & 0.744 & 0.855 & 0.773 \\
Playability / chords     & 0.855 & 0.859 & 0.856 & 0.816 & 0.882 & 0.865 & 0.942 & 0.893 \\
Setup / action           & 0.809 & 0.830 & 0.805 & 0.817 & 0.838 & 0.839 & 0.912 & 0.903 \\
Shipping damage          & 0.887 & 0.874 & 0.899 & 0.889 & 0.913 & 0.888 & 0.929 & 0.912 \\
String quality           & 0.894 & 0.898 & 0.891 & 0.869 & 0.906 & 0.905 & 0.948 & 0.907 \\
Tuning stability         & 0.924 & 0.935 & 0.950 & 0.921 & 0.950 & 0.933 & 0.974 & 0.944 \\
Visual appearance        & 0.845 & 0.847 & 0.845 & 0.801 & 0.867 & 0.838 & 0.913 & 0.872 \\
\bottomrule
\end{tabular}
\end{table*}

\section{ABSA vs.\ LLM Topic Sentiment Comparison}
\label{app:absa}

Table~\ref{tab:absa} compares the two topic-level sentiment approaches.
ABSA scores are systematically inflated by the dataset's positive skew;
the two methods agree at the extremes but diverge in the middle.

\begin{table}[h]
\centering
\small
\setlength{\tabcolsep}{3pt}
\caption{Topic positive\% under ABSA (BART zero-shot) vs.\
LLM few-shot (Mistral-7B). Sorted by LLM positive\%.}
\label{tab:absa}
\begin{tabular}{lcc}
\toprule
\textbf{Topic} & \textbf{ABSA Pos\%} & \textbf{LLM Pos\%} \\
\midrule
Acoustic tone            & 80.6 & 96.2 \\
Beginner learning        & 79.2 & 91.3 \\
Pickups                  & 76.9 & 88.8 \\
Guitar size              & 79.2 & 82.0 \\
Playability / chords     & 79.6 & 81.4 \\
Setup / action           & 78.4 & 79.9 \\
Visual appearance        & 78.5 & 68.9 \\
Accessories              & 70.6 & 67.7 \\
Electronics / controls   & 77.4 & 62.8 \\
Customer service / ret.\ & 70.7 & 47.3 \\
Fret / neck setup        & 76.1 & 43.5 \\
Tuning stability         & 73.2 & 33.2 \\
String quality           & 74.6 & 27.6 \\
Shipping damage          & 65.6 & 17.4 \\
\bottomrule
\end{tabular}
\end{table}

\section{Replication: Per-Topic Sentiment}
\label{app:rep_topics}

Tables~\ref{tab:topics_drums} and~\ref{tab:topics_keyboards} report
the full topic-level positive/negative distributions for the drums
and keyboards corpora, scored with the fine-tuned RoBERTa
classifier on the LDA-assigned reviews.
Functional-quality topics (\textit{Product defects/returns},
\textit{Returns/broken keys}, \textit{Durability/stopped working})
again sit at the bottom of the ranking, mirroring the
\textit{Shipping damage} / \textit{String quality} / \textit{Tuning
stability} cluster in guitars.

\begin{table}[h]
\centering
\small
\setlength{\tabcolsep}{4pt}
\caption{Topic-level sentiment, drums.}
\label{tab:topics_drums}
\begin{tabular}{lrcc}
\toprule
\textbf{Topic} & \textbf{N} & \textbf{Pos\%} & \textbf{Neg\%} \\
\midrule
Gift / holiday purchase           & 23{,}389 & 97.0 &  3.0 \\
Beginner / ease of play           & 12{,}469 & 92.1 &  7.9 \\
Singing bowls / meditation        & 15{,}278 & 91.3 &  8.7 \\
Drum hardware / shells            & 10{,}335 & 88.1 & 11.9 \\
Cymbals                           & 10{,}587 & 87.6 & 12.4 \\
Small percussion accessories      & 15{,}261 & 85.5 & 14.5 \\
Electronic drums / practice pad   &  9{,}753 & 82.2 & 17.8 \\
Product defects / returns         & 16{,}646 & 44.6 & 55.4 \\
\bottomrule
\end{tabular}
\end{table}

\begin{table}[h]
\centering
\small
\setlength{\tabcolsep}{4pt}
\caption{Topic-level sentiment, keyboards.}
\label{tab:topics_keyboards}
\begin{tabular}{lrcc}
\toprule
\textbf{Topic} & \textbf{N} & \textbf{Pos\%} & \textbf{Neg\%} \\
\midrule
Gift / holiday purchase           & 6{,}100 & 99.1 &  0.9 \\
Kids music lessons                & 2{,}559 & 97.1 &  2.9 \\
Synthesizer / sound design        & 2{,}267 & 92.3 &  7.7 \\
Beginner learning                 & 7{,}742 & 90.8 &  9.2 \\
MIDI controller / live perform.\  & 1{,}577 & 89.7 & 10.3 \\
Key feel / weighted keys          & 4{,}979 & 87.8 & 12.2 \\
Brand quality / acoustic feel     & 4{,}814 & 83.8 & 16.2 \\
Sustain pedal                     &   826   & 82.4 & 17.6 \\
Sound quality / speaker           & 3{,}746 & 80.1 & 19.9 \\
Appearance / color / display      & 1{,}293 & 78.3 & 21.7 \\
Connectivity (USB/MIDI/phone)     & 1{,}676 & 68.4 & 31.6 \\
Power supply \& accessories       & 3{,}581 & 66.8 & 33.2 \\
Durability / stopped working      & 1{,}741 & 57.3 & 42.7 \\
Returns / broken keys             & 5{,}696 & 31.4 & 68.6 \\
\bottomrule
\end{tabular}
\end{table}

\section{Pointwise Rubric and Two-Step Prompt}
\label{app:rubric}

\paragraph{Pointwise rubric (1--5, judge anchored).}
\textbf{Emotional Resonance} (ELM peripheral / Rossiter--Percy):
1 = no emotional cue, 3 = mild/generic affect, 5 = vivid and
on-strategy.
\textbf{Argument Strength} (ELM central): 1 = unsupported claim,
3 = mixed specific/generic, 5 = concrete causal reasoning.
\textbf{Specificity}: 1 = could describe any product, 3 = balanced,
5 = product-specific terms throughout.
\textbf{Purchase Motivation} (AIDA desire/action): 3 is
deliberately calibrated as ``neutral'' to avoid optimistic drift on
a 2--3 sentence ad format.

\paragraph{Two-step ad generation (per product$\times$topic).}
\textit{Step 1 (analysis).} The generator reads up to 10 positive
and 10 negative real reviews for the topic and outputs
\texttt{insight\_a} (most compelling genuine praise) and
\texttt{insight\_b} (most common real concern with a feasible
counter).
\textit{Step 2 (generation).} A single call returns Strategy~A and
Strategy~B together, with both insights as context, forcing active
contrast and preventing convergence.
Step 1 is the chain-of-thought scaffold and is grounded in actual
review evidence rather than free generation.

\section{Product$\times$Topic Selection Rule}
\label{app:select}

For every subcategory we (i)~enumerate all
product$\times$topic pairs with $\geq$5 positive
\textit{and} $\geq$5 negative reviews; (ii)~rank products by
\texttt{score = n\_qualifying\_topics + description\_length / 1000};
(iii)~deduplicate by brand keeping the highest-scoring product;
(iv)~retain the top three products with at least 2 qualifying
topics and a non-empty metadata description; (v)~for each retained
product pick the (up to) 3 qualifying topics minimising
$|\text{pct\_positive}-50\%|$ to maximise input parity between
strategies.
This is fully automated (no manual curation), so the procedure
ports unchanged across guitars, drums, and keyboards.

\section{Reproducibility and Resources}
\label{app:repro}

\paragraph{Splits.}
All sentiment classifiers use stratified 80/20 splits with
\texttt{random\_state=42}.
The default test set holds 24{,}815 reviews
(3{,}698 negative / 21{,}117 positive).
Fine-tuned RoBERTa and Mistral-7B use the same split applied to the
24{,}849 rating-3-free rows, which differs by 34 samples
($<$0.15\%) and does not affect conclusions.

\paragraph{LDA hyperparameters.}
LDA is trained with Gensim using $\alpha{=}$asymmetric,
$\eta{=}$auto, 10~passes, evaluated with the C$_v$ coherence metric
\cite{roder2015coherence}.
The grid search covers
$K \in \{5, 10, 15, 20, 25\}$ and four $(\alpha, \eta)$
configurations.
Topic assignment uses a probability threshold of 0.15.

\paragraph{Classical baseline tuning.}
TF-IDF + LR best parameters: $C{=}10$, max\_features=100k,
ngram\_range=(1,2), max\_df=0.85.
Sentence Transformer embeddings use the
\texttt{all-mpnet-base-v2} model.

\paragraph{Fine-tuned RoBERTa.}
RoBERTa-base, 3 epochs, learning rate $2{\times}10^{-5}$, linear
warmup, early stopping on validation Macro~F1.

\paragraph{Mistral 5-shot.}
Mistral-7B-Instruct-v0.3 with temperature~0 and a fixed 5-shot
exemplar set drawn from the training split.

\paragraph{ABSA threshold.}
BART-large-MNLI assigns a topic label only when entailment
probability is at least 0.70.
Without this threshold, low-confidence predictions inflate noise.

\paragraph{Per-topic error rates (RoBERTa).}
Hardest topics: Shipping damage 7.1\%, Customer service / returns
6.8\%.
Easiest topics: Acoustic tone 1.0\%, Setup / action 1.3\%.

\paragraph{Stage-2 API setup.}
Generator and judge models in Stage~2 are accessed through their
public APIs at the time of writing.
Ad order is randomised and each pair is judged twice
(A-first and B-first) to control for position bias
\cite{zheng2023judging}.

\end{document}
